% ****************************************************************************************************
% Capítulo 9: Para Quem a Medicação Trabalha?
% ****************************************************************************************************
\chapter{Para Quem a Medicação Trabalha?}
\label{ch:para-quem}

% ****************************************************************************************************
\section{A Pergunta Não Formulada}
\label{sec:pergunta-nao-formulada}
% ****************************************************************************************************

Há uma pergunta que deveria preceder toda prescrição psiquiátrica, mas que raramente é formulada: \textit{para quem esta medicação trabalha?} A resposta parece óbvia --- para o paciente, naturalmente. Mas a obviedade, aqui como em outros lugares, merece suspeita.

Considere: um idoso com demência em instituição de longa permanência recebe um antipsicótico para \enquote{agitação}. A agitação cessa. Objetivo terapêutico alcançado? Talvez. Mas a quem beneficia, exatamente, que esse idoso não esteja mais agitado? A ele --- que agora está sedado, com risco aumentado de queda, pneumonia, morte? Ou à instituição --- que agora não precisa lidar com comportamento desafiador, não precisa investigar a causa da agitação, não precisa alocar funcionários para cuidado individualizado?

Este capítulo examina duas dimensões interligadas: a \textit{banalização} de certas medicações --- particularmente antipsicóticos --- para usos muito além de suas indicações originais; e a \textit{pressuposição} de que a medicação é etapa em um processo de cuidado, quando frequentemente é substituto desse processo.

% ****************************************************************************************************
\section{A Banalização do Antipsicótico}
\label{sec:banalizacao}
% ****************************************************************************************************

\subsection{Quetiapina: De Antipsicótico a \enquote{Sonífero de Luxo}}

A quetiapina (Seroquel) foi aprovada para tratamento de esquizofrenia, transtorno bipolar e, como adjuvante, depressão maior. São condições graves, que afetam profundamente a vida de quem as tem. O antipsicótico, nesse contexto, é ferramenta potente para situações potentes.

E no entanto.

Em estudo conduzido em lares de idosos na Suíça, pesquisadores encontraram que 94,3\% das prescrições de quetiapina eram \foreign{off-label} --- ou seja, para indicações não aprovadas. As principais: agitação, demência, ansiedade, depressão, insônia. Apenas 4,8\% das prescrições correspondiam a indicações para as quais a droga foi desenvolvida e testada.

A quetiapina tornou-se o que alguns chamam, com ironia amarga, de \enquote{sonífero de luxo} --- prescrita para insônia sem evidência robusta de eficácia, sem aprovação regulatória para essa indicação, mas com todos os riscos metabólicos e cardiovasculares de um antipsicótico.

\subsection{A Contenção Química em Instituições}

Se a prescrição ambulatorial de quetiapina para insônia é questionável, seu uso institucional para \enquote{manejo comportamental} é francamente alarmante.

Análise de mais de 12.000 lares de idosos nos Estados Unidos encontrou que mais de um em cada cinco residentes estava recebendo antipsicótico. A taxa de prescrição era quase dez vezes maior que a taxa de diagnósticos para os quais esses medicamentos são considerados apropriados.

O termo técnico para esse uso é \textit{contenção química} --- paralelo farmacológico da contenção física (amarras, grades) que a medicina moderna aprendeu a considerar violação de direitos. A relatora especial da ONU para direitos de pessoas idosas, Claudia Mahler, em relatório de 2024, recomendou a proibição de contenções químicas em instituições de cuidado.

\subsection{Para Quem É Conveniente Que o Idoso Durma?}

A pergunta que não se faz: \textit{por que} esse idoso está agitado? Está com dor não tratada? Está com infecção urinária? Está assustado em ambiente desconhecido? Está tentando comunicar necessidade que não consegue verbalizar?

Um ensaio clínico em lares de idosos na Noruega demonstrou que abordagem escalonada para tratamento de dor pode reduzir sintomas comportamentais em residentes com demência. Tratar a dor resolve o \enquote{comportamento problema}. Mas tratar dor exige avaliação, atenção, tempo. O antipsicótico é mais rápido.

Há aqui uma inversão perversa. O \enquote{comportamento problema} --- agitação, agressividade, vagar noturno --- é, frequentemente, a única forma que resta a uma pessoa com demência para comunicar que algo está errado. Silenciar essa comunicação com sedativo não resolve o problema; apenas torna invisível quem o tem.

% ****************************************************************************************************
\section{O Controle Disfarçado de Tratamento}
\label{sec:controle-tratamento}
% ****************************************************************************************************

\subsection{A Lógica do \enquote{Manejável}}

Há uma palavra que aparece com frequência perturbadora em discussões sobre uso de psicotrópicos em populações vulneráveis: \textit{manejável}. O paciente se tornou mais \enquote{manejável}. O comportamento agora é \enquote{manejável}.

A pergunta que a palavra evita: \textit{manejável por quem?}

A mesma lógica aparece em outros contextos. Antipsicóticos são amplamente prescritos para crianças e adultos com deficiência intelectual ou autismo, frequentemente para \enquote{comportamentos desafiadores}. O NHS britânico, em sua campanha STOMP, documentou que 17,1\% dos adultos com deficiência intelectual recebiam antipsicóticos --- a grande maioria sem diagnóstico de condição psicótica.

Comportamentos classificados como \enquote{desafiadores} frequentemente são respostas compreensíveis a ambientes inadequados, dor não reconhecida, comunicação não atendida. Medicar o comportamento sem abordar sua causa é silenciar o mensageiro em vez de ouvir a mensagem.

\subsection{A Diferença Entre Sedar e Tratar}

Há uma distinção fundamental que a prática contemporânea frequentemente obscurece: a diferença entre \textit{sedar} e \textit{tratar}.

Tratar implica abordar a causa do sofrimento. Sedar implica suprimir sua expressão. São operações radicalmente distintas, com implicações éticas igualmente distintas.

Quando um idoso com demência recebe antipsicótico para \enquote{agitação}, o que exatamente está sendo tratado? A demência não melhora. A causa da agitação não é investigada. O que acontece é sedação: redução global de atividade, incluindo a atividade que estava sendo rotulada como \enquote{problema}.

Isso não é tratamento em nenhum sentido rigoroso da palavra. É controle comportamental através de química cerebral. É, para usar termo que a medicina contemporânea prefere evitar, \textit{contenção}.

\subsection{O Risco Que Não Entra na Conversa}

Antipsicóticos em idosos com demência aumentam risco de morte. Isso não é hipótese; é fato estabelecido por múltiplos estudos, reconhecido por agências regulatórias, motivo de \foreign{black box warnings} nas bulas. A magnitude do risco --- aumento de 60--70\% na mortalidade --- é substancial.

E no entanto, essa informação frequentemente não chega à família que consente com o uso da medicação. A conversa, quando existe, tende a ser sobre \enquote{ajudar com a agitação}, \enquote{melhorar a qualidade de vida}. Não sobre aumento de risco de pneumonia, de eventos cardiovasculares, de morte.

% ****************************************************************************************************
\section{A Janela Que Ninguém Atravessa}
\label{sec:janela-vazia}
% ****************************************************************************************************

\subsection{O Modelo Implícito}

Há um modelo implícito que sustenta grande parte da prática psicofarmacológica contemporânea:

\begin{enumerate}
\item O paciente está em crise
\item A medicação estabiliza (reduz sintomas agudos)
\item Com a estabilização, abre-se \enquote{janela terapêutica}
\item Através dessa janela, a psicoterapia promove transformação duradoura
\item O paciente melhora de forma sustentável
\end{enumerate}

O modelo é elegante. O problema é que depende de uma premissa que frequentemente não se verifica: que haverá alguém do outro lado da janela.

\subsection{O Problema do Passo 3}

Quem fará o trabalho psicoterapêutico? Na teoria, o mesmo profissional que prescreveu, ou um colega em trabalho colaborativo. Na prática, há uma cisão crescente entre quem medica e quem faz psicoterapia.

Estudo publicado em 2025 documenta tendência importante: nos Estados Unidos, a proporção de pacientes que recebiam psicoterapia de psiquiatras caiu de 41\% para 34\% entre 2018 e 2021.

No sistema público brasileiro, acesso a psicólogo é frequentemente restrito. No sistema privado, cobertura de psicoterapia é frequentemente limitada.

Resultado: muitos pacientes recebem a medicação, mas não a psicoterapia. A janela terapêutica se abre --- e permanece vazia.

\subsection{A Medicação Como Substituto}

Entre adultos em tratamento ambulatorial de saúde mental nos Estados Unidos, quase dois terços recebiam apenas medicação, sem psicoterapia. A psicoterapia --- o trabalho relacional, psicológico, existencial que deveria dar sentido à \enquote{janela terapêutica} --- chega a fração pequena da população que precisa dela.

O que acontece quando a medicação funciona (reduz sintomas), mas não há ninguém para atravessar a janela que ela abre? O paciente melhora sintomaticamente, mas não há trabalho sobre os padrões de pensamento, sobre os conflitos relacionais, sobre o contexto existencial. A medicação se torna não facilitador de transformação, mas substituto dela.

% ****************************************************************************************************
\section{Uma Farmacologia a Serviço do Sujeito}
\label{sec:servico-sujeito}
% ****************************************************************************************************

\subsection{Critérios para Prescrição Centrada no Paciente}

Proponho critérios para o que chamarei de \textit{prescrição centrada no paciente}:

\textbf{Primeiro critério: Clareza sobre indicação.} Antes de prescrever, perguntar: esta medicação está sendo prescrita para tratar condição do paciente, ou para manejar situação que incomoda outros?

\textbf{Segundo critério: Investigação de causas.} Antes de medicar comportamento, investigar o que o comportamento comunica. Agitação pode ser dor. Agressividade pode ser medo. Medicar sem investigar é silenciar mensageiro antes de ouvir mensagem.

\textbf{Terceiro critério: Continuidade de cuidado.} Antes de abrir janela terapêutica, verificar se haverá alguém para atravessá-la.

\textbf{Quarto critério: Consentimento informado real.} Informar não apenas benefícios esperados, mas riscos reais, incluindo os graves. Permitir participação genuína na decisão.

\textbf{Quinto critério: Reavaliação contínua.} Perguntar periodicamente: esta medicação ainda está servindo ao paciente? Ou está apenas sendo continuada por inércia?

\subsection{A Pergunta Que Deveria Preceder Toda Receita}

O capítulo começou com uma pergunta: \textit{para quem esta medicação trabalha?}

A resposta honesta, em muitos casos, é desconfortável. A medicação frequentemente trabalha para o sistema --- para a instituição que precisa de corpos dóceis, para o psiquiatra que precisa de consultas rápidas, para o plano de saúde que prefere pílulas a sessões.

Isso não significa que medicação seja sempre inapropriada. Significa que a pergunta deve ser feita --- explicitamente, conscientemente, a cada prescrição. E que a resposta deve ser, genuinamente, que a medicação serve ao paciente. Não ao paciente em abstrato, como categoria diagnóstica a ser tratada, mas ao paciente concreto, como pessoa com história, valores, contexto, e direito de participar das decisões sobre sua própria vida.

% ****************************************************************************************************
\section{A Medicação e Seus Senhores}
\label{sec:medicacao-senhores}
% ****************************************************************************************************

A farmacologia psiquiátrica é conquista notável. Drogas que há meio século não existiam permitem hoje que pessoas com condições graves vivam vidas que antes eram impossíveis. Reconhecer isso não é capitular ao farmacologismo; é reconhecer realidade.

Mas reconhecer conquistas não impede crítica. A crítica desenvolvida neste capítulo não é antifarmacológica. É, se quisermos usar termo, \textit{farmacologia crítica}: reflexão sobre os usos da farmacologia, sobre os valores que incorpora, sobre os interesses que serve.

A conclusão não é \enquote{não prescreva}. É \enquote{pergunte, antes de prescrever}. Pergunte pelo tempo do paciente, não apenas pelo tempo da molécula. Pergunte pelo preço que o paciente está disposto a pagar. Pergunte para quem a medicação trabalha --- e só prescreva quando a resposta for, honestamente, \enquote{para o paciente}.

Há, nessa exigência, algo de retorno ao fundamento hipocrático: \foreign{primum non nocere}, primeiro não prejudicar. Mas também algo mais: reconhecimento de que o paciente é sujeito de sua vida, não objeto de intervenção. Que tem valores próprios, tempos próprios, prioridades próprias. E que a medicina --- incluindo a psiquiatria, incluindo a farmacologia --- existe para servir a esses valores, tempos e prioridades. Não o contrário.

% ****************************************************************************************************
% Referências do Capítulo
% ****************************************************************************************************
\vspace{2em}
\begin{small}
\noindent\textsc{Referências}

\vspace{1em}
\noindent CMS - Centers for Medicare \& Medicaid Services. (2025). Updates to nursing home care compare.

\noindent Husebo, B. S., et al. (2011). Efficacy of treating pain to reduce behavioural disturbances in nursing homes with dementia. \textit{BMJ}, 343, d4065.

\noindent Mahler, C. (2024). Rights of older persons. Report of the Independent Expert. United Nations Human Rights Council.

\noindent Müller, R., et al. (2020). Off-label use of quetiapina in nursing homes. \textit{British Journal of Clinical Pharmacology}, 86(6), 1165--1172.

\noindent Olfson, M., et al. (2025). Trends in outpatient psychotherapy: 2018--2021. \textit{American Journal of Psychiatry}.

\noindent Schneider, L. S., et al. (2005). Risk of death with atypical antipsychotic drug treatment for dementia. \textit{JAMA}, 294(15), 1934--1943.
\end{small}
